% New research fragment. The general weighted statement remains open.
\section{Weighted terminal paths and short-cycle averaging}
\label{sec:path-edge-half-weighted}

Let $N$ be a connected simple graph with three degree-two terminals.
Every other vertex has degree three.
Write $c=c(N)$, and let $p_{ij}$ count simple paths between terminals $i,j$.
For an edge $e$, use $c_e$ and $p_{ij}(e)$ for the corresponding counts through $e$.
The proposed weighted inequality is
\begin{equation}\label{eq:path-edge-half-target}
 2c_e+2\sum_{i<j}w_{ij}p_{ij}(e)
 \le c+1+\sum_{i<j}w_{ij}p_{ij},\qquad w_{ij}\ge0.
\end{equation}
The edge may depend on the weights.
The statements below do not prove this inequality for arbitrary $N$.

\subsection{An equivalent finite convex problem}

\begin{lemma}\label{lem:path-edge-half-dual}
Inequality~\eqref{eq:path-edge-half-target} holds for every nonnegative weight triple if and only if there is a probability distribution $(\alpha_e)_{e\in E(N)}$ such that
\[
 \sum_e\alpha_ec_e\le\frac{c+1}{2},\qquad
 \sum_e\alpha_ep_{ij}(e)\le\frac{p_{ij}}2
 \quad(1\le i<j\le3).
\]
\end{lemma}

\begin{proof}
Set
\[
 v_e=\bigl(2c_e-c-1,
 2p_{12}(e)-p_{12},2p_{13}(e)-p_{13},2p_{23}(e)-p_{23}\bigr).
\]
The desired distribution exists exactly when the convex hull of the $v_e$ meets the nonpositive orthant.
If these sets are disjoint, strict separation gives a nonzero vector $\lambda\ge0$ with $\lambda\cdot v_e>0$ for every edge.
If $\lambda_0=0$, add a sufficiently small positive amount to that coordinate.
The inequalities remain strict because there are finitely many edges.
After dividing by $\lambda_0$, the remaining coordinates give weights that violate~\eqref{eq:path-edge-half-target}.
Conversely, the displayed distribution makes the weighted mean of the left side minus the right side nonpositive.
At least one edge therefore satisfies the required inequality.
\end{proof}

\subsection{Exact identities on a triangle or square}

For a vertex set $U$, let $p_{ij}^{\neg U}$ count terminal paths disjoint from $U$.
This count is zero if $U$ contains either endpoint.

\begin{proposition}\label{prop:path-edge-half-short-rim}
Suppose $W$ is a cycle of length $q\in\{3,4\}$ in $N$.
Then
\begin{align}
 \sum_{e\in E(W)}c_e
 &=\frac q2\bigl(c+1-c(N-V(W))\bigr),\label{eq:path-edge-half-rim-cycles}\\
 \sum_{e\in E(W)}p_{ij}(e)
 &=\frac q2\bigl(p_{ij}-p_{ij}^{\neg V(W)}\bigr)
 \quad(1\le i<j\le3).\label{eq:path-edge-half-rim-paths}
\end{align}
Consequently, the uniform distribution on $E(W)$ satisfies Lemma~\ref{lem:path-edge-half-dual}.
In particular, the weighted inequality holds for all nonnegative weights.
\end{proposition}

\begin{proof}
First suppose that $W$ has a chord.
Then $q=4$, and the chord forms a triangle in $N$.
For the conclusion about weighted inequalities, this triangle suffices.
For the identities, however, retain the original square and regard its chord as an outside edge.

Fix the edges outside $E(W)$ of a terminal path.
Their parity, together with the prescribed terminal pair, determines an even subset $J\subseteq V(W)$.
The selected rim edges have odd-degree set $J$ within the rim.
Every rim vertex has at most one outside incidence.
A terminal on the rim has none.

If $|J|=2$, the two rim solutions are complementary arcs.
Both complete the fixed outside edges to a simple terminal path.
Indeed, the interior vertices of either arc have no selected outside incidence or terminal condition.
The two paths use $q$ rim edges in total.
If $|J|=4$, then $q=4$, and each feasible rim solution is a two-edge matching.
One or both solutions may give a single path.
Each valid solution still uses exactly $q/2$ rim edges.
If $J$ is empty, a terminal path uses no rim edge.
It then avoids every rim vertex, by the degree restrictions above.
These cases prove~\eqref{eq:path-edge-half-rim-paths}.

Apply the same partition to cycles.
The cases $|J|=2$ and $|J|=4$ again have mean rim usage $q/2$.
For $J$ empty, the only cycle using rim edges is $W$ itself.
Every other cycle in this case avoids $V(W)$.
The cycle $W$ contributes $q$ rim edges.
This proves~\eqref{eq:path-edge-half-rim-cycles}.
\end{proof}

The identities give the more precise weighted conclusion
\[
 \min_{e\in E(W)}\left(2c_e+2\sum_{i<j}w_{ij}p_{ij}(e)\right)
 \le c+1+\sum_{i<j}w_{ij}p_{ij}
       -c(N-V(W))-\sum_{i<j}w_{ij}p_{ij}^{\neg V(W)}.
\]
Thus the unresolved weighted case has girth at least five.

\subsection{Scope of the proposed inequality}

Let $G$ be simple and cubic, and suppose $G-v$ is connected.
Take $N=G-v$, with terminals the three neighbours of $v$.
Every cycle through $v$ corresponds to exactly one terminal path in $N$.
With all three weights equal to one,~\eqref{eq:path-edge-half-target} would therefore imply
\[
 \min_{e\not\ni v}2c_e(G)\le c(G)+1.
\]
This would prove the edge-half inequality whenever $c(G)$ is even.
For odd $c(G)$, it gives the integer upper bound $(c(G)+1)/2$.
The weighted statement therefore extends beyond graphs with nontrivial three-edge cuts.

For positive weights, define
\[
 \lambda_1=\sqrt{w_{12}w_{13}/w_{23}},\quad
 \lambda_2=\sqrt{w_{12}w_{23}/w_{13}},\quad
 \lambda_3=\sqrt{w_{13}w_{23}/w_{12}}.
\]
Then $w_{ij}=\lambda_i\lambda_j$.
Adding a new vertex adjacent to the terminals interprets the weighted expression as a cycle count with three star-edge weights.
All internal edges retain weight one.
This reformulation supplies no additional restriction on a positive weight triple.

\subsection{The exact pentagon obstruction}

Suppose now that $W$ is an induced five-cycle.
Partition cycles, and each terminal-path family, by their selected edges outside $E(W)$.
For each such outside edge set, there are at most two rim completions.
When the required odd-degree set on the rim has four vertices, one completion has two edges and the other has three.
Call these the short and long completions, respectively.

Let $L_c,S_c$ count outside edge sets with exactly one valid cycle completion, which is long or short, respectively.
For terminal pair $i,j$, define $L_{ij},S_{ij}$ in the same way, requiring a single terminal path.
Only the four-vertex odd-degree cases enter these counts.
Then
\begin{align}
 \sum_{e\in E(W)}c_e
 &=\frac52\bigl(c+1-c(N-V(W))\bigr)+\frac12(L_c-S_c),
 \label{eq:path-edge-half-pentagon-cycles}\\
 \sum_{e\in E(W)}p_{ij}(e)
 &=\frac52\bigl(p_{ij}-p_{ij}^{\neg V(W)}\bigr)
       +\frac12(L_{ij}-S_{ij}).
 \label{eq:path-edge-half-pentagon-paths}
\end{align}
Indeed, two valid complementary completions have mean rim length $5/2$.
A single short or long completion differs from that mean by $-1/2$ or $1/2$.
The zero- and two-vertex odd-degree cases behave exactly as in the short-rim proof.

Consequently, uniform averaging on this pentagon proves the weighted inequality if and only if
\begin{equation}\label{eq:path-edge-half-pentagon-target}
 (L_c-S_c)+\sum_{i<j}w_{ij}(L_{ij}-S_{ij})
 \le5\left(c(N-V(W))+\sum_{i<j}w_{ij}p_{ij}^{\neg V(W)}\right).
\end{equation}
This is an exact criterion for that averaging distribution.
It is not a necessary condition for the existence of a good individual edge elsewhere.

The defect terms have a concrete switching interpretation.
A long-only terminal path becomes one terminal path and one vertex-disjoint cycle after symmetric difference with $E(W)$.
The resulting pair uses precisely the two short rim edges.
Conversely, every such pair whose complementary rim completion is connected supplies one long-only path.
The short-only case has the analogous three-edge rim description.
For cycles, the invalid completion consists of two vertex-disjoint cycles.
These descriptions are bijections because symmetric difference with $E(W)$ is an involution.
They do not yet compare the positive and negative defect counts.

\subsection{The terminal-local criterion and actual shore restrictions}

For a terminal $i$, let $c_i$ count cycles containing that vertex.
Let $v_i$ count paths between the other terminals $j,k$ that visit $i$.
Put
\[
 A_i=2c_i-c-1,\qquad B_i=2v_i-p_{jk}.
\]
Average the two incident-edge coefficient vectors in Lemma~\ref{lem:path-edge-half-dual}.
Each path ending at $i$ uses exactly one of these edges.
Each path between $j,k$ uses both or neither.
Each cycle uses both or neither.
The averaged weighted surplus is exactly
\[
 A_i+w_{jk}B_i.
\]
Thus $A_i\le0$ and $B_i\le0$ give one distribution that works for every weight triple.
The girth-five Petersen subdivision example in the companion computational notes disproves this sufficient condition in general.

When a pole is one shore of a three-edge cut in a three-edge-connected cubic graph, its capped graph is three-edge-connected.
Equivalently, for every nonempty proper vertex subset $U$ of the pole, with $t$ terminals in $U$,
\[
 |\delta_N(U)|\ge\max\{t,3-t\}.
\]
In particular, the pole has no bridge, and every two-edge cut separates one terminal from the other two.
If the capped graph also has girth at least five, the pole has girth at least five and its terminals have pairwise distance at least three.
These are genuine structural restrictions on opposite shores.
The general weighted statement does not assume them.
